% PDF pp. 19--24; continuation of proof of 5.1.
Thus $K$ is of finite type over $S$, hence of finite presentation over $S$, since $S$ is locally noetherian. Consequently, by EGA IV$_3$, 8.8.2.4, to show that there exists an open neighborhood $U$ of $s$ such that $K_U=H_U$, it suffices to show that $K_{S'}=H_{S'}$, where $S'$ is the spectrum of $A=\OO_{S,s}$. One may therefore suppose $S$ local with closed point $s$. Replacing the noetherian local ring $A$ by its completion $\widehat A$ if necessary, which introduces a faithfully flat and quasi-compact base change $\widehat S\to S$, one may even suppose\nde{For $K=H$ if $K_{\widehat S}=H_{\widehat S}$, cf. SGA 1, VIII 5.4; but what follows does not use the hypothesis ``$A$ complete''.} $A$ complete if one wishes.

By 3.4 one has $u_n=v_n$ for every $n$, where, as usual, the subscript $n$ indicates reduction modulo $\goth m^{n+1}$ ($\goth m$ being the maximal ideal of $A$). For every integer $m>0$, denoting by ${}_mu,{}_mv$ the homomorphisms ${}_mH\to G$ induced by $u,v$, one therefore also has $({}_mu)_n=({}_mv)_n$. This being true for every $n$, and ${}_mH$ being finite over $S$ by 2.2, it follows that ${}_mu={}_mv$ by 5.0. This being true for every $m$, one therefore has $u=v$ by 4.7.

b) $G$ of finite presentation.\nde{The original has been expanded in what follows.} Since $H$ is also of finite presentation over $S$, by EGA IV$_3$, 8.8.2, we may again suppose $S$ local with closed point $s$ and prove that then $u=v$. If $f:S'\to S$ is a faithfully flat quasi-compact morphism, and if one denotes by $f_H$ the morphism $H'\to H$ and by $u',v':H'\to G'$ the morphisms deduced from $u,v$, then the equality $u'=v'$ implies $u'\circ f_H=v'\circ f_H$, whence $u=v$, since $f_H$ is an epimorphism. Thus, making a faithfully flat and quasi-compact extension of the base if necessary, one may suppose in addition $H$ diagonalizable, hence of the form $D_S(M)$, with $M$ a commutative group of finite type.
% typo?: u' composed with f_H retained as printed.
Introduce, as in the proof of 4.6, the increasing filtering family of $\ZZ$-subalgebras of finite type $A_i$ of $A=\OO_{S,s}$, and $S_i=\Spec(A_i)$.\footnotemark[\value{footnote}] Note that $H=D_S(M)$ comes, for every $i$, from the diagonalizable group $H_i=D_{S_i}(M)$. Since $G$ is of finite presentation over $S$, by EGA IV$_3$, 8.8.2 (see also VI$_{\mathrm B}$, 10.2 and 10.3), there exists an index $i$, a group prescheme $G_i$ of finite presentation over $S_i$, and morphisms of $S_i$-groups $u_i,v_i:H_i\to G_i$ from which $u,v$ come by base change. Let $s_i$ be the image of $s$ in $S_i$, and let $\rho_i:H_i\times_{S_i}G_i\to G_i$ be the morphism of $S_i$-preschemes defined by $\rho_i(h,g)=u_i(h)g u_i(h)^{-1}$. Then, since $u_s$ is central, $\rho_s=\rho_i\times_{\kappa(s_i)}\kappa(s)$ equals the second projection, and hence the same is true of $\rho_i$ (since $\kappa(s_i)\to\kappa(s)$ is faithfully flat and quasi-compact), i.e. $(u_i)_{s_i}$ is central. Similarly, since $u_s=v_s$, one has $(u_i)_{s_i}=(v_i)_{s_i}$. One may then apply a) to the situation over $S_i$, whence immediately the announced conclusion.
% typo?: rho_i rather than its fiber retained as printed.
\begin{corollary}{5.2}\label{IX.5.2}
Let $u:H\to G$ be a homomorphism of $S$-group preschemes, with $H$ of multiplicative type and of finite type, and $S$ locally noetherian or $G$ of finite presentation over $S$. Let $s\in S$ and suppose that $u_s$ is the unit homomorphism. Then there exists an open neighborhood $U$ of $s$ such that $u_U$ is the unit homomorphism.
\end{corollary}
\begin{corollary}{5.3}\label{IX.5.3}
Let $u,H,G$ be as in 5.2, and suppose in addition $G$ separated over $S$. Let $U$ be the set of $s\in S$ such that $u_s$ is the unit homomorphism. Then $U$ is an open and closed subset of $S$, and $u_U:H_U\to G_U$ is the unit homomorphism.
\end{corollary}
It remains only to prove that $U$ is closed. Now let $K=\Ker(u)$; since $G$ is separated over $S$, $K$ is a closed subprescheme of $H$, and $U$ is the set of $s\in S$ such that $K_s=H_s$. One then easily sees that $U$ is closed, for example as an application of VIII 6.4 ($H$ being essentially free over $S$ by VIII 6.3), or by noting that one may suppose $S$ reduced and $U$ dense in $S$, hence schematically dense in $S$, which implies $H_U$ schematically dense in $H$ since $H$ is flat over $S$,\nde{And of finite presentation over $S$, cf. EGA IV$_3$, 11.10.5 (ii) and 11.9.10 (ii).} and as $u$ and the unit homomorphism coincide on $H_U$, they coincide on $H$.
\begin{corollary}{5.4}\label{IX.5.4}
Let $S$ be a prescheme, $H$ and $G$ two $S$-groups, with $H$ of multiplicative type and of finite type and $G$ separated and of finite presentation over $S$, $\pi:S'\to S$ a\nde{``Covering morphism for the faithfully flat and quasi-compact topology'' has been replaced by ``universal effective epimorphism'', in order to be able to apply this to the morphism $K\to S$ of 5.5, and the beginning of the proof has been modified accordingly.} universal effective epimorphism (for example a faithfully flat and quasi-compact morphism, or a morphism admitting a section, cf. IV 1.12) with geometrically connected fibers.\nde{This is the case, for example, if $S=\Spec k$ and $S'=\Spec k'$, where $k$ is a field and $k'$ a radicial extension of $k$, cf. X, Prop. 1.4.}

Let $u':H'\to G'$ be a central homomorphism of the $S'$-groups deduced from $H,G$ by the base change $S'\to S$. Then there exists a unique homomorphism $u:H\to G$ of $S$-groups such that $u\times_S S'=u'$. When $S'/S$ admits a section $g$, $u$ is the morphism deduced from $u'$ by the base change $g:S\to S'$.
\end{corollary}
Since $\pi$ is a universal effective epimorphism, the same is true of $H'\to H$, whence uniqueness of $u$, cf. the beginning of the proof of 5.1 b). If $\pi$ admits a section $g$, then $u'=\pi^*(u)$ implies $u=g^*\pi^*(u)=g^*(u')$.

For the existence of $u$, one is reduced, by IV 2.3, to showing that the two homomorphisms $u''_1,u''_2:H''\to G''$ of $S''$-groups deduced from $u'$ by the two base changes $\operatorname{pr}_1,\operatorname{pr}_2:S''=S'\times_S S'\to S'$ are identical. Now they coincide on the diagonal of $S''$; more precisely, the inverse images of $u''_1$ and $u''_2$ by the diagonal morphism $S'\to S''$ are identical (being equal to $u'$).\nde{The following sentence has been added.} Since $u''_1$ and $u''_2$ are central, one may apply 5.3 to the morphism $u''_1(u''_2)^{-1}$. There therefore exists an open and closed subset $U$ of $S''$, containing the diagonal of $S''$, such that $u''_1$ and $u''_2$ coincide over $U$.

Now the fibers of $S'/S$ being geometrically connected, the same is true of those of $S''/S$, which are therefore a fortiori connected, whence it follows that $U$ (containing the diagonals of the said fibers) contains the said fibers, and hence equals $S''$, which completes the proof.
\begin{corollary}{5.5}\label{IX.5.5}
Let $S$ be a prescheme, $K$ an $S$-group prescheme locally of finite type\nde{The hypothesis ``locally of finite type'' has been added to agree with the reference VI$_{\mathrm A}$, 2.4. In fact, one can show that over a field every connected group scheme is geometrically connected.} and with connected fibers, $H$ an invariant subgroup of multiplicative type and of finite type of $K$. Then $H$ is a central subgroup of $K$.
\end{corollary}
\nde{The original has been expanded in what follows.} Let us first note that $\pi:K\to S$ has geometrically connected fibers, since for a group scheme locally of finite type over a field, connected implies geometrically connected (cf. VI$_{\mathrm A}$ 2.4). One may then apply 5.4 by putting $G=H$ and $S'=K$ to the homomorphism of $K$-groups $u':H_K\to H_K$ defined set-theoretically by $(h,k)\mto(khk^{-1},k)$, which is central since $H$ is commutative. The inverse image of $u'$ by the unit section $\varepsilon:S\to K$ is the identity homomorphism of $K$, so by 5.4 the same is true of $H_K\to H_K$, hence $H$ is central in $K$.
% typo?: "identity homomorphism of K" retained as printed.
Let us state the variants of the preceding results for central subgroups of multiplicative type. One obtains, proceeding as for the preceding results (and using 3.2 bis):
\begin{theorem}{5.1 bis}\label{IX.5.1bis}
Let $G$ be an $S$-group prescheme, $H_1$ and $H_2$ two subgroup preschemes of multiplicative type and of finite type, with $(H_1)_s$ central. Suppose $S$ locally noetherian or $G$ of finite presentation over $S$. Then for every $s\in S$ such that $(H_1)_s=(H_2)_s$ there exists an open neighborhood $U$ of $s$ such that $H_1|U=H_2|U$.
\end{theorem}
\begin{corollary}{5.3 bis}\label{IX.5.3bis}
Under the preceding conditions, suppose in addition $G$ separated over $S$. Let $U$ be the set of $s\in S$ such that $(H_1)_s=(H_2)_s$. Then $U$ is an open and closed subset of $S$, and $H_1|U=H_2|U$.
\end{corollary}
\begin{corollary}{5.4 bis}\label{IX.5.4bis}
Let $S$ be a prescheme, $G$ a group prescheme separated and of finite presentation over $S$, $S'\to S$ a covering morphism for the faithfully flat quasi-compact topology, with geometrically connected fibers, $H'$ a subgroup of multiplicative type and of finite type of $G'=G\times_S S'$.

Then there exists a unique subgroup $H$ of $G$ such that $H'=H\times_S S'$, and $H$ is of multiplicative type and of finite type over $S$, by 1.1 and 2.1 b). If $S'\to S$ admits a section $g$, $H$ is the inverse image of the subgroup $H'$ of $G'$ by $g:S\to S'$.
\end{corollary}
\begin{theorem}{5.6}\label{IX.5.6}
Let $S$ be a prescheme, $u:H\to G$ a homomorphism of $S$-group preschemes, with $H$ of multiplicative type and of finite type, $G$ of finite presentation over $S$.

a) Suppose $G$ to have connected fibers, and let $s\in S$ be such that $u_s:H_s\to G_s$ is a central homomorphism. Then there exists an open neighborhood $U$ of $s$ such that the homomorphism $H|U\to G|U$ induced by $u$ is central.

b) Suppose that for every $s\in S$, $u_s:H_s\to G_s$ is central. Then $u$ is central.
\end{theorem}
\nde{The original has been expanded in the proof of a).} Let us prove (a). Proceeding exactly as in the proof of 5.1 (b), one may suppose $S$ local with closed point $s$, then $H=D_S(M)$ diagonalizable, then that there exists a $\ZZ$-subalgebra of finite type $A_i$ of $A=\OO_{S,s}$ and a morphism $u_i:D_{S_i}(M)\to G_i$ such that $u_{s_i}$ is central (where $s_i$ is the image of $s$ in $S_i$) and from which $u$ comes by base change. This reduces us to the case where $S$ is local noetherian with closed point $s$, and one must then prove that $u$ is central.

Let $K$ be the subprescheme $\Ker(v,w)$, where $v,w:H\times_S G\rightrightarrows G$ are defined by
\[
v(h,g)=g,\qquad w(h,g)=\operatorname{int}(u(h))\cdot g=u(h)g u(h)^{-1}.
\]
Then $K$ is a sub-$G$-group of the $G$-group $H_G=H\times_S G$, and one wants to show that it equals $H_G$ itself. By 4.9 one is reduced to proving that it contains the ${}_m(H_G)=({}_mH)\times_S G$ for every integer $m>0$, which reduces us to the case where $H={}_mH$, hence $H$ finite over $S$.

Let $e$ be the unit element of the fiber $G_s$; then $S'=\Spec(\OO_{G,e})$ is a noetherian local scheme ($G$ being of finite presentation over noetherian $S$); put $S'_n=S'\times_S S_n$, where $S_n=\Spec(A/\goth m^{n+1})$. Then $K_{S'}=K\times_G S'$ is a subprescheme of $H_{S'}=H\times_S S'$, and by 3.9 one has $K_{S'_n}=H_{S'_n}$ for every $n$.\nde{The original has been expanded and corrected, replacing ``unit section of $(G_H)_n$ over $H_n$'' by $H_{S'_n}$.} Since $H_{S'}$ is finite over $S'$, one concludes from 5.0 (applied to the noetherian local ring $B$) that $K_{S'}=H_{S'}$.

On the other hand, since $H\times_S G$ is noetherian (being of finite presentation over noetherian $S$), the immersion $K\hookrightarrow H\times_S G$ is of finite type (cf. EGA I, 6.3.5), so $K$ is of finite type, hence of finite presentation over $G$. Then the equality $K_{S'}=H_G\times_G S'$ implies, by EGA IV$_3$, 8.8.2.4, that there exists an open neighborhood $W$ of $e$ in $G$ such that $K\times_G W=H_G\times_G W=H\times_S W$. Thus $K$ contains the open neighborhood $V=H\times_S W$ of the unit section of $G_H$ over $H$.

For every $t\in H$, the fiber $G_t$ (being an algebraic $\kappa(t)$-group) is Cohen-Macaulay (VI$_{\mathrm A}$, 1.1.1), and hence without embedded components; as it is in addition connected, hence irreducible (VI$_{\mathrm A}$, 2.4), it has its generic point as its unique associated point. Thus by EGA IV$_2$, 3.1.8, the open subset $V_t$ is schematically dense in $G_t$. By 4.3, $V$ is therefore schematically dense in $G_H$. Moreover, since $K$ is a sub-$H$-group of $G_H$, it induces on each fiber $G_h$ a subgroup $K_h$, and as the latter contains an open neighborhood of the neutral element and $G_h$ is connected, it follows that $K_h=G_h$, whence $K=G_H$ set-theoretically. Thus $K$ is a closed subprescheme of $G_H$ containing the schematically dense subprescheme $V$, hence $K=G_H$, which proves a).

Finally, b) is a direct consequence of 3.9, in view of the fact that to check that a subprescheme $K$ of $P=G\times_S H$ is identical to $P$, it suffices to check that for every base change $S'\to S$, where $S'$ is an artinian quotient of a local ring of $S$, one has $K'=P'$.

\sgasection{6}{Monomorphisms of groups of multiplicative type, and canonical factorization of a homomorphism of such a group}
\begin{lemma}{6.1}\label{IX.6.1}
Let $S$ be a quasi-compact prescheme, $G$ a commutative $S$-group prescheme, of finite presentation and quasi-finite over $S$. Then there exists an integer $n>0$ such that $n\cdot\mathrm{id}_G=0$, i.e. $G={}_nG$.
\end{lemma}
If $S$ is the spectrum of a field $k$, $G$ is finite over $k$, and by VII$_{\mathrm A}$ 8.5 it suffices to take $n=\deg(G/k)$. Suppose now $S=\Spec(A)$ local artinian; let $k$ be the residue field of $A$, $G_0=G\otimes_A k$, and $n_0=\deg(G_0/k)$. Distinguish two cases.

a) $k$ is of characteristic zero. Then $G_0$ is separable over $k$, hence $G$ unramified over $S$. Then the unit section of $G$ is an open immersion, hence ${}_nG=\Ker(n\cdot\mathrm{id}_G)$ is an open subscheme of $G$, so for it to equal $G$ it suffices that it do so set-theoretically; one may therefore then take $n=n_0$.

b) $k$ is of characteristic $p>0$. Denote by $\goth m$ the maximal ideal of $A$ and put $S_r=\Spec(A/\goth m^{r+1})$ for every $r$. Let $m$ be an integer such that $\goth m^{m+1}=0$; I say that one may take $n=p^m n_0$. Indeed, proceeding by induction on $m$ and putting $n'=p^{m-1}n_0$, one may suppose that $n'\cdot\mathrm{id}_G$ induces the zero endomorphism in $G_{m-1}=G\times_S S_{m-1}$.

\nde{The original has been expanded in what follows, specifying the results of Exp. III that are used.} Denote by $E$ the set of endomorphisms $u$ of $G$ having this property and by $K$ the kernel of the morphism of group functors $G\to\prod_{S_{m-1}/S}G$ (cf. III 0.1.2). Then $E$ is identified with $\Hom_{S\text{-}\mathrm{gr.}}(G,K)$; in particular the abelian group law on $E$ is induced by that of $K$. Now by III 0.9, $K$ is the $S$-group functor associating with every $g:T\to S$ the $k$-vector space
\[
\Hom_{\OO_{T_0}}(g_0^*(\Omega^1_{G_0/S_0}),\goth m^m\OO_T);
\]
one therefore has $pu=0$ for every $u\in E$. Thus one has $pn'\cdot\mathrm{id}_G=0$, i.e. $p^m n_0\cdot\mathrm{id}_G=0$.

Suppose now $S$ noetherian (one reduces to this in 6.1 by the usual reduction to the noetherian case\nde{Since $S$ is quasi-compact, it is covered by finitely many affine open subsets, so one is reduced to the case where $S=\Spec(A)$. Then, since $G$ is of finite presentation over $S$, one may apply EGA IV$_3$, 8.8.2.}). It then suffices to combine the preceding with the
\begin{lemma}{6.2}\label{IX.6.2}
Let $X$ be a prescheme quasi-finite over noetherian $S$, and consider an increasing filtering family of subpreschemes $(X^i)_{i\in I}$ having the following property:
\[
\begin{gathered}
\text{for every }s\in S\text{ and }n\geq0,\text{ putting }S_{s,n}=\Spec(\OO_{S,s}/\goth m^{n+1}),\text{ there exists }i\in I\\
\text{such that }X^i\times_S S_{s,n}=X\times_S S_{s,n}.
\end{gathered}
\]
Under these conditions there exists an $i\in I$ such that $X^i=X$.
\end{lemma}
Since $S$ is noetherian, there exists a maximal open subset $U$ such that one has $X|U=X^i|U$ for large $i$; one shall show that $U=S$. In other words, one shall show that if $U\ne S$, one can find a $U'$ strictly larger than $U$, and an $i$ such that $X|U'=X^i|U'$. Localizing at a maximal point $s$ of $S-U$, one is reduced to the case where $S$ is local with closed point $s$ and $U=S-\{s\}$. (Indeed, if one denotes by $S'=\Spec(\OO_{S,s})$ and if there exists $i$ such that $X\times_S S'=X^i\times_S S'$, then\nde{$X-\{s\}$ has been corrected above to $S-\{s\}$. On the other hand, recall that a quasi-finite morphism is supposed of finite type, cf. EGA II, 6.2.3 (and EGA III$_2$, Err$_{\mathrm{III}}$ 20 for the definition of ``locally quasi-finite''). Thus here ($S$ being noetherian), $X$ is of finite presentation over $S$, each immersion $X^i\hookrightarrow X$ is of finite type (by EGA I, 6.3.5), so $X^i$ is of finite presentation over $S$, and one may apply EGA IV$_3$, 8.8.2.} there exists an open neighborhood $V$ of $s$ such that $X|V=X^i|V$; thus, taking $i$ sufficiently large that $X^i|U=X|U$, one shall have $X|W=X^i|W$, where $W=U\cup V$.)

Then for large $i$, since $X|U=X^i|U$, one sees that $X^i$ is a closed subprescheme of $X$ defined by an ideal $\cal J^{(i)}$ with support contained in $X_s=X_0=X\times_S S_0$ (where $S_n=\Spec(A/\goth m^{n+1})$, $S=\Spec(A)$). As $X_0$ is quasi-finite over $S_0$, $X_0$ is a finite closed subset of noetherian $X$, so $\cal J^{(i)}$ is a module of finite length. It follows that there exists an integer $n\geq0$ such that $\cal J^{(i)}\cap\goth m^{n+1}\OX=0$. On the other hand, by the hypothesis in 6.2 one may suppose (on increasing $i$) that the image of $\cal J^{(i)}$ in $\OO_{X_n}=\OX/\goth m^{n+1}\OX$ is zero. This implies $\cal J^{(i)}=0$, hence $X^i=X$. \hfill Q.E.D.
\begin{lemma}{6.3}\label{IX.6.3}
Let $S$ be a prescheme, $K$ a group prescheme over $S$, locally of finite presentation, $s\in S$ such that $K_s$ is quasi-finite (resp. unramified) over $\kappa(s)$ at the unit element. Then there exists an open neighborhood $U$ of $s$ such that $K|U$ is locally quasi-finite (resp. unramified) over $U$.\footnote{Cf. VI$_{\mathrm B}$ 2.5 (ii) for a more general statement.}
\end{lemma}
Indeed, let $V$ be the set of points $x$ of $K$ such that, denoting by $t$ the image of $x$ in $S$, the fiber $K_t$ is quasi-finite (resp. unramified) over $\kappa(t)$ at $x$, i.e. such that $x$ is isolated in $K_t$ (resp. and its local ring in $K_t$ a separable extension of $\kappa(t)$). One knows that $V$ is open since $K$ is locally of finite presentation over $S$,\nde{Cf. EGA IV$_3$, 13.1.4 and EGA IV$_4$, 17.4.1. Moreover, the image of $x$ has been denoted by $t$ (instead of $s$).} so if $\varepsilon$ denotes the unit section of $K$, $\varepsilon^{-1}(V)$ is open. By hypothesis it contains $s$, so it is an open neighborhood $U$ of $s$. The latter does the job; in other words $t\in U$ implies that $K_t$ is locally quasi-finite (resp. unramified) over $\kappa(t)$: indeed, since $K_t$ is a group locally of finite type over $\kappa(t)$, this follows from the fact that it is quasi-finite (resp. unramified) over $\kappa(t)$ at the point $\varepsilon(t)$, cf. VI$_{\mathrm B}$ 1.3.

Combining 6.1 and 6.3 one finds the
\begin{theorem}{6.4}\label{IX.6.4}
Let $S$ be a prescheme, $H$ an $S$-group of multiplicative type and of finite type, $K$ a closed subgroup prescheme, of finite presentation over $S$. Let $s\in S$ be such that $K_s$ is finite.

Then there exists an open neighborhood $U$ of $s$ such that $K|U$ is contained in ${}_nH|U$ for some $n$, and a fortiori (${}_nH$ being finite over $S$) such that $K|U$ is finite over $U$.
\end{theorem}
Using Nakayama's lemma, one deduces:
\begin{corollary}{6.5}\label{IX.6.5}
With the preceding notation, suppose that $K_s$ is the unit group. Then there exists an open neighborhood $U$ of $s$ such that $K|U$ is the unit group.
\end{corollary}
\begin{corollary}{6.6}\label{IX.6.6}
Let $u:H\to G$ be a homomorphism of $S$-group preschemes, with $H$ of multiplicative type and of finite type, and $G$ separated over $S$. Suppose in addition $S$ locally noetherian, or $G$ locally of finite presentation\nde{In fact it suffices to suppose $G$ locally of finite type over $S$; by EGA IV$_4$, 1.4.3 (v), this implies ($H\to S$ being of finite presentation) that $H\to G$ is locally of finite presentation, and hence also $K\to S$, which is deduced from it by base change.} over $S$.
% Statement continues on PDF p. 25, in en-05.tex.
